% ============================================================================
\section{第一批 syscall handler}
\label{sec:syscall-handlers}
% ============================================================================

本阶段实现三个 handler 的初始版本。它们在 D-2 阶段是\textbf{占位实现}（stub），
功能有限但足以验证 dispatch 路径正确。D-3（\texttt{int 0x80} 后端）完成后，
这些 handler 会被真正的 Ring 3 代码所调用。

% ----------------------------------------------------------------------------
\subsection{\texttt{sys\_exit}：终止进程}
% ----------------------------------------------------------------------------

\codefile{src/kernel/core/syscall.c}
\begin{lstlisting}[style=cstyle]
/*
 * sys_exit -- 终止当前进程
 *
 * arg1 (EBX) = exit status
 *
 * 当前没有进程管理，exit 只能 halt。
 * 将来 D-7/D-8 会做 ZOMBIE 状态转换和资源回收。
 */
static int32_t sys_exit(const syscall_regs_t* regs) {
    printk("[syscall] exit(status=%d)\n", (int)regs->arg1);

    for (;;) {
        __asm__ volatile("hlt");
    }

    return 0;  /* unreachable */
}
\end{lstlisting}

\texttt{sys\_exit} 目前直接 halt——因为还没有进程管理，没有"下一个进程"可以调度。
将来实现 PCB 和调度器后，\texttt{exit} 会将进程状态设为 ZOMBIE 并释放资源。

% ----------------------------------------------------------------------------
\subsection{\texttt{sys\_write}：写数据}
% ----------------------------------------------------------------------------

\begin{lstlisting}[style=cstyle]
/*
 * sys_write -- 写数据到文件描述符
 *
 * arg1 (EBX) = fd        (当前只支持 1=stdout, 2=stderr)
 * arg2 (ECX) = buf       (用户态缓冲区地址)
 * arg3 (EDX) = count     (字节数)
 *
 * 返回值 = 实际写入字节数，-1=错误
 */
static int32_t sys_write(const syscall_regs_t* regs) {
    uint32_t fd    = regs->arg1;
    uint32_t buf   = regs->arg2;
    uint32_t count = regs->arg3;

    /* 只接受 stdout(1) 和 stderr(2) */
    if (fd != 1 && fd != 2) {
        return -1;
    }

    /* 安全检查：buf 必须来自用户空间 */
    if (buf >= 0xC0000000u) {
        printk("[syscall] write: bad user pointer 0x%08x\n", buf);
        return -1;
    }

    /* 上限截断，防止恶意大 count */
    if (count > 4096) {
        count = 4096;
    }

    const char* p = (const char*)(uintptr_t)buf;
    for (uint32_t i = 0; i < count; i++) {
        printk("%c", p[i]);
    }

    return (int32_t)count;
}
\end{lstlisting}

\begin{warnbox}
\textbf{用户指针安全检查}是系统调用中最重要的安全措施之一。

\texttt{buf >= 0xC0000000} 是粗略的检查——内核空间从 \hex{C0000000} 开始，
用户态缓冲区地址不应该落在这个范围内。如果恶意用户程序传入一个内核地址，
\texttt{sys\_write} 会拒绝而不是泄露内核数据。

将来实现 VFS 后，这里还需要更精细的检查：
验证 \texttt{[buf, buf+count)} 整个范围都在合法的用户 VMA 中。
\end{warnbox}

% ----------------------------------------------------------------------------
\subsection{\texttt{sys\_brk}：堆管理}
% ----------------------------------------------------------------------------

\begin{lstlisting}[style=cstyle]
/*
 * sys_brk -- 调整进程数据段末尾
 *
 * arg1 (EBX) = new_brk (0 = 查询当前 brk)
 * 返回值 = 当前 brk 地址
 *
 * 当前只返回占位值。将来 F-6（mmap）会和 VMA 集成。
 */
static int32_t sys_brk(const syscall_regs_t* regs) {
    (void)regs;
    printk("[syscall] brk(0x%08x) -- stub\n", regs->arg1);
    return 0x00400000;
}
\end{lstlisting}

\texttt{brk} 是 Unix 传统的堆管理方式——移动数据段末尾来分配/释放内存。
用户态的 \texttt{malloc} 内部就是调用 \texttt{brk}（小分配）或 \texttt{mmap}（大分配）。
